\subsection{Fontes Públicas Vigentes}

O CVE Program constitui a base de identificação do ecossistema, pois fornece o identificador comum usado para referenciar vulnerabilidades publicamente divulgadas \cite{cve_program_overview}. Seu papel é funcionar como âncora semântica entre diferentes fontes. Entretanto, o registro CVE, isoladamente, não resolve todas as necessidades analíticas, pois informações como severidade, fraqueza associada, produto afetado, exploração conhecida, probabilidade de exploração e versões corrigidas são distribuídas entre outras bases.

A National Vulnerability Database (NVD), mantida pelo NIST, complementa os registros CVE com informações técnicas padronizadas, incluindo métricas CVSS, categorias CWE, configurações CPE, referências e histórico de mudanças \cite{nvd_home,first_cvss,mitre_cwe,nvd_cpe_dictionary}. Por isso, ela é frequentemente usada como núcleo técnico de estudos quantitativos sobre vulnerabilidades. Ainda assim, a NVD não deve ser tratada como fonte única. O próprio CVSS é uma medida de severidade, não de risco, e mudanças recentes na operação do NVD reforçam a necessidade de complementar seus dados com outras fontes \cite{nvd_cvss_metrics,nist_nvd_operations_2026}. A partir de 15 de abril de 2026, o NIST passou a priorizar o enriquecimento de CVEs presentes no catálogo KEV da CISA, de softwares usados pelo governo federal dos Estados Unidos ou de softwares considerados críticos \cite{nist_nvd_operations_2026}. Essa priorização tende a aumentar a assimetria de completude entre registros.

O GitHub Advisory Database e o OSV ampliam a visão centrada em CVE/CPE ao representar vulnerabilidades no nível de ecossistemas de software, pacotes e faixas de versão \cite{github_advisory_database_docs,github_advisory_database_repo,osv_database,osv_schema}. Essa granularidade é especialmente importante para estudos sobre software open source, cadeias de dependência e análise de versões vulneráveis ou corrigidas. Enquanto o GitHub Advisory Database agrega advisories com GHSA IDs, CVEs, CWEs, pacotes, ecossistemas, versões vulneráveis, versões corrigidas e severidade \cite{github_advisory_database_docs}, o OSV propõe um formato estruturado para mapear vulnerabilidades diretamente a versões de pacotes ou commits \cite{osv_schema}.

A CISA Known Exploited Vulnerabilities (KEV) acrescenta uma dimensão operacional distinta: a presença de uma vulnerabilidade no catálogo indica evidência de exploração conhecida \cite{cisa_kev_catalog}. Essa informação é valiosa porque permite diferenciar vulnerabilidades apenas severas daquelas que já possuem exploração observada. No entanto, a ausência de um CVE no KEV não significa ausência de risco; significa apenas que a vulnerabilidade não foi incluída naquele catálogo.

O EPSS, por sua vez, fornece um score probabilístico que estima a chance de exploração de uma vulnerabilidade em curto prazo \cite{first_epss}. Sua contribuição é complementar ao CVSS e ao KEV: enquanto o CVSS descreve severidade técnica, e o KEV registra exploração conhecida, o EPSS estima probabilidade de exploração \cite{first_cvss,nvd_cvss_metrics,cisa_kev_catalog,first_epss}.

Outras fontes, como Exploit-DB, Vulners e VirusTotal, têm papéis diferentes. O Exploit-DB é relevante para identificar disponibilidade pública de exploits ou provas de conceito, mas não comprova exploração real em ambiente operacional \cite{exploitdb}. Vulners funciona como plataforma agregadora de inteligência de vulnerabilidades, mas possui restrições de acesso, API e licenciamento que reduzem sua adequação como fonte primária para um dataset aberto \cite{vulners,vulners_api_docs}. VirusTotal é útil para enriquecimento por IOCs, reputação e relações entre arquivos, URLs, domínios e IPs, mas não é uma base primária de CVEs \cite{virustotal_api_v3}.

\begin{table}[htbp]
    \centering
    \caption{Comparação entre fontes públicas}
    \label{tab:fontes-publicas}
    \small
    \renewcommand{\arraystretch}{1.15}
    \resizebox{\textwidth}{!}{%
        \begin{tabular}{|p{2.6cm}|p{2.5cm}|p{2.4cm}|p{3.2cm}|p{3.3cm}|p{3.3cm}|}
            \hline
            \textbf{Fonte}            & \textbf{Escopo principal}                                     & \textbf{Identificadores}      & \textbf{Pontos fortes}                                                     & \textbf{Limitações}                                                                                     \\
            \hline
            CVE Program / CVE List V5 & Registro canônico de vulnerabilidades publicamente divulgadas & CVE-ID                        & Autoridade do identificador, padronização e ampla adoção                   & Pouco enriquecimento técnico em comparação com NVD, EPSS, KEV e bases de advisories                     \\
            \hline
            NVD                       & Enriquecimento técnico padronizado                            & CVE, CVSS, CWE, CPE           & Severidade, fraquezas, produtos afetados, configurações e referências      & CVSS não mede risco; completude pode variar; priorização seletiva a partir de 2026                      \\
            \hline
            GitHub Advisory Database  & Advisories de segurança, principalmente open source           & GHSA, CVE, CWE                & Pacote, ecossistema, versões vulneráveis/corrigidas e contexto de advisory & Cobertura concentrada em ecossistemas suportados; exige filtragem de malware e advisories não revisados \\
            \hline
            OSV                       & Vulnerabilidades open source por pacote, versão e commit      & OSV ID, CVE, GHSA             & Boa semântica de versões e integração multi-ecossistema                    & Menos adequado para produtos generalistas fora do modelo de pacotes                                     \\
            \hline
            CISA KEV                  & Vulnerabilidades com exploração conhecida                     & CVE                           & Sinal operacional forte de exploração confirmada                           & Cobertura seletiva e binária                                                                            \\
            \hline
            EPSS                      & Probabilidade de exploração                                   & CVE                           & Score probabilístico e histórico diário                                    & Mede ameaça/probabilidade, não risco completo                                                           \\
            \hline
            Exploit-DB                & Exploits públicos e provas de conceito                        & EDB-ID, CVE quando disponível & Evidência de exploit público                                               & PoC não equivale a exploração real; cobertura incompleta                                                \\
            \hline
            Vulners                   & Inteligência agregada de vulnerabilidades                     & IDs internos, CVE, CPE        & Correlação multi-fonte                                                     & Restrições de API, plano e licença                                                                      \\
            \hline
            VirusTotal                & Threat intelligence e IOCs                                    & Hash, URL, domínio, IP        & Reputação e relações entre artefatos                                       & Não é CVE-first; API pública limitada                                                                   \\
            \hline
        \end{tabular}%
    }

    \vspace{0.2cm}
    {\footnotesize\textit{Fonte: elaboração própria com base nas documentações oficiais das fontes analisadas \cite{cve_program_overview,nvd_home,github_advisory_database_docs,github_advisory_database_repo,osv_database,osv_schema,cisa_kev_catalog,first_epss,exploitdb,vulners,vulners_api_docs,virustotal_api_v3}.}}

\end{table}

Essa comparação evidencia que nenhuma fonte, isoladamente, cobre de forma satisfatória todas as dimensões necessárias para pesquisa acadêmica em vulnerabilidades, veja na Tabela \ref{tab:fontes-publicas}. O CVE fornece identidade; o NVD fornece severidade e enriquecimento técnico; o GitHub Advisory Database e o OSV fornecem granularidade por pacote; o KEV indica exploração conhecida; o EPSS estima probabilidade de exploração; e fontes como Exploit-DB, Vulners e VirusTotal adicionam sinais contextuais \cite{cve_program_overview,nvd_home,github_advisory_database_docs,osv_database,osv_schema,cisa_kev_catalog,first_epss,exploitdb,vulners,virustotal_api_v3}.
